\ssscp{Medial *NC clusters}{13-02-01-02}
\ili{As} mentioned above,
the treatment of clusters where the first consonant is a nasal
is not included in Blust’s claims for PCEMP,
but we provide a discussion of these here to round off our discussion
of medial CC clusters.

Throughout the Wallacean region there are sporadic retentions
of medial *NC → NC,
with the nasal and consonant often being modified slightly from the historical forms.
For example, in \ili{Kambera},
the medial *NC clusters are retained with post-nasal voicing
as in example \qf{13-11}.
The data are from \citet[54]{Blust2008}.

\noindent{\wg\st{6pt}\tblr{>{\hspace{-\tabcolsep}}l<{\hspace{-\tabcolsep}}
>{\hspace{-\tabcolsep}\enspace}
lll}
\lsptoprule
%\tbx{13-11}	&	PMP	&	\ili{Kambera}	&	gloss	\\	\midrule
%\endfirsthead
\tbx{13-11}	&	PMP	&	\ili{Kambera}	&	gloss	\\	\midrule
%\endhead
	&	*kə\tbr{mp}uŋ	&	ka\tbr{mb}u	&	belly	\\
	&	*tu\tbr{mb}uq	&	tu\tbr{mb}u	&	grow	\\
	&	*pa\tbr{nd}ak	&	pa\tbr{nd}ak-u	&	short	\\
	&	*sə\tbr{nd}iR	&	hɐ\tbr{nd}i	&	lean against 	\\
	&	*di\tbr{ŋd}iŋ	&	di\tbr{nd}iŋ-u	&	wall	\\
	&	*pa\tbr{ŋd}an	&	pɐ\tbr{nd}aŋ-u	&	pandanus	\\
	&	*bə\tbr{ŋk}uq	&	bɐ\tbr{ŋg}u	&	bent, curved	\\
	&	*na\tbr{ŋk}a	&	na\tbr{ŋg}a	&	jackfruit	\\
\lspbottomrule
\end{tabular}\mbox{}\\}

The evidence from \ili{Kambera} is sufficient to affirm Blust’s assessment
that the medial clusters need to be reconstructed for the region
when a nasal is involved.
This pattern is found in several Wallacean subgroups.
In previous chapters, along with the summary in \srf{13-01-02},
we showed that reflexes for medial *-mp-
and *-mb- are often different from the reflexes of *-p-
and *-b- respectively, and need to be kept distinct.
Similarly for *-nt- and *-t-, *-nd-,
and *-d-, *-ns- and *-s-, *-nz-
and *-z-, *-nj- and *-j-, *-ŋk-
and *-k-, and *-ŋg- and *-g-.
Many of these latter clusters are extremely rare,
so finding clear reflexes is a challenge,
particularly for making comparative claims across languages or subgroups.
Many etyma however, have doublets *(N)C (like *upu
and *umpu; *ŋisi and *ŋinsi), with both forms needed to account
for the data in the region,
and sometimes in neighbouring languages.

But the medial *NC clusters are also more complicated than the
\ili{Kambera} example described above,
even when reduction is involved.
For example, data from the \tsc{Seram-Tanimbar-Bomberai} language
\ili{Kei} Kecil shows that a single language may have more than one
reflex of a medial *NC cluster in the same etymon,
as in the following example,
from \citet{Grimes1991a}: PMP *ma-di\tbb{ŋ}\tbr{d}iŋ ‘cold’ >
\it{ŋa-ri\tbr{d}in} ‘cold’, \it{na-b-ri\tbb{n}in} ‘chills’.

\ili{Buru} also shows that a single sound correspondence cannot be
assumed for similar medial *NC sequences in the same language.
\ili{Buru} has no root internal medial CC clusters,
so it reduces medial clusters,
even in loans, as in \ili{Malay} \it{ganti} > \ili{Buru} \it{gati} ‘replace,
change’; \ili{Malay} \it{tuŋgu} ‘wait,
keep watch, guard’ > \ili{Buru} \it{tugu} ‘guard (against intrusion), keep watch, wait’.
Phonotactic pressure means that in medial historical *CC of any
combination, one of the consonants will have to be lost,
or the two will have to merge.
The regular sound correspondence for \ili{Buru} is PMP *d > \it{r},
based on data in \citet{GrimesGrimes2020} in which initial
*d- > \it{r} (22 instances),
medial *-d- > \it{r} (25 instances),
and final *-d > {\0} (18 instances).
Yet note the following examples with the retention of medial
*d = \it{d} in some reduplicated monosyllables in example \qf{13-12},
but not in others.

\noindent{\wg\st{6pt}\begin{tabularx}
{\linewidth}{>{\hspace{-\tabcolsep}}l<{\hspace{-\tabcolsep}}
>{\hspace{-\tabcolsep}\enspace}lll
>{\raggedright\arraybackslash}X}
\lsptoprule
%\tbx{13-12}	&	PMP gloss	&	PMP	&	\ili{Buru}	&		\\	\midrule
%\endfirsthead
\tbx{13-12}	&	PMP gloss	&	PMP	&	\ili{Buru}	&		\\	\midrule
%\endhead
	&	cold	&	*ma-di\tbr{ŋd}iŋ	&	eb-ri\tbr{d}i(n)	&	cold (without \it{-n} = V; with \it{-n} = N)	\\
	&	gloom, darkness	&	*də\tbr{md}əm	&	re\tbr{d}e(n)	&	darkness, gloom	\\
	&	warm self near fire	&	*da\tbr{nd}an/*da\tbr{ŋd}aŋ	&	ra\tbb{r}a	&	get warm over fire (undergoer)	\\
	&		&		&	eʔ-ra\tbb{r}a	&	warm oneself over fire (actor)	\\
	&	fever, illness	&	**da\tbr{md}am (PWMP)	&	ra\tbb{r}a(n)	&	fever, heat	\\
\lspbottomrule
\end{tabularx}\mbox{}\\}

The first two etyma in example \qf{13-12} that retain *d = \it{d}
rather than the normal *d > \it{r} illustrate what \citet[293f]{Wolff2010}
describes as the medial *C being “protected by a nasal” such
that it does not undergo its normal expected sound change.
Like \ili{Kambera}, the \ili{Buru} data also show that the *NC clusters must
be reconstructed for languages in the region.
Both parts of the *NC combination are required to explain the
surface forms, even though there is reduction.
And like \ili{Kei} Kecil,
the \ili{Buru} data also show that a single language may have more
than one correspondence for these medial clusters,
so they need to be evaluated on an item-by-item basis,
since so many exceptions to broader claims appear sporadically
throughout the region.

Even the idea that medial *NC clusters in reduplicated monosyllables
in the region “assimilated to the place of articulation of the
obstruent” \citep[246]{Blust1993} needs to be treated with a bit of caution.
First of all, such assimilation is also widespread west
and north of the Blust line,
as in PMP *diŋdiŋ ‘wall of a house,
partition off’ > Kapampangan \it{dindiŋ},
Tombonuwo \it{rindiŋ}, Murut (Paluan) \it{rindiŋ}, Supan \it{rindiŋ},
\ili{Ngaju} \it{dindiŋ}, \ili{Malagasy} \it{rindrina},
Iban \it{dindiŋ}, \ili{Malay} \it{dindiŋ},
\ili{Sasak} \it{dindiŋ} (data from \citealp{BlustTrussel2020}).
Furthermore, within the target region it appears the nasal may
occasionally be retained unassimilated,
as in the Southwest \ili{Tarangan}
and \ili{Dobel} forms for ‘dark’ in \trf{13-08},
although these two could possibly be from an unreduplicated root
*dəm.\footnote{
		In \citet{BlustTrussel2020}, the root *-dəm\sub{1}
		‘dark; overcast’ is only listed with preceding
		material, quite different from the Southwest \ili{Tarangan}
		and \ili{Dobel} forms.}
The collective data together show a variety
of patterns, including coalescence (→ C\sub{3}) in the Kaitetu
and \ili{Saparua} forms in what happens to the medial cluster,
and loss or haplology in \ili{Tii} and \ili{Galolen}.
(Data in \trf{13-08} are adapted,
updated and expanded from \citet{Grimes1991a}; there are a range of
local glosses all relating to the historical glosses,
such as ‘cold, chills,
cool to the touch’,
and ‘thunder, rumbling sound’).
Reconstructing the medial cluster as a cluster better accounts
for the collective data in the region than does reconstructing
some sort of reduction that only accounts for part of the data.

\begin{table}\tcp{Reflexes of *-CC- clusters in Wallacea}{13-08}
%\begin{adjustbox}{width=\textwidth}
	\begin{threeparttable}%\st{2pt}
	\begin{tabular}{ll|lll@{ }l}\lsptoprule
	&	gloss	&	cold	&	dark, gloom	&	thunder	&	\\
	&	PMP	&	*ma-di\tbr{ŋd}iŋ	&	*də\tbr{md}əm	&	*du\tbr{Rd}uR\su{†} 	&	\\ \midrule
\tsc{BL}	&	\ili{Kambera}	&	‹mari\tbr{ng}u›	&	{\hr}ru\tbr{nd}uŋ-u	&	{\hr}ndu\tbr{r}u	&	\\
\tsc{BL}	&	\ili{Dhao}	&	{\hr}mari\tbr{ŋ}i	&		&	{\hr}ɗo\tbr{r}o	&	\\
\tsc{TB}	&	\ili{Tii}	&		&		&	{\hr}ruu	&	\\
\tsc{TB}	&	\ili{Helong}	&	{\hr}b-li\tbr{ŋ}in	&	{\hr}de\tbr{d}eŋ	&	{\hr}lo\tbr{t}o	&	\\
\tsc{TB}	&	\ili{Tetun}	&	{\hr}ma-li\tbr{r}in	&		&	{\hr}taru\tbr{t}u	&	\\
\tsc{TB}	&	\ili{Galolen}	&	{\hr}na-diin	&		&		&	\\
\tsc{TB}	&	\ili{Rasua}	&	{\hr}na-ri\tbr{ɲ}i	&		&		&	\\
\tsc{TB}	&	\ili{Wetan}	&	{\hr}ri\tbr{nn}i	&		&		&	\\
\tsc{TB}	&	\ili{Selaru}	&	{\hr}mrin	&		&	{\hr}kdu\tbr{d}u,	&	{\hr}‘rumble’\\
	&		&		&		&	{\hr}kru\tbr{t}u	&	{\hr}‘thunder’\\
\tsc{Aru}	&	SW \ili{Tarangan}	&		&	{\hr}de\tbr{m}ur	&	{\hr}du\tbr{d}im-na	&	\\
\tsc{Aru}	&	\ili{Dobel}	&		&	{\hr}de\tbr{m}ur	&	{\hr}fadu\tbr{d}um,	&	{\hr}‘thunder’\\
	&		&		&		&	{\hr}du\tbr{d}um	&	{\hr}‘sound’\\
\tsc{Aru}	&	\ili{Kola}	&		&	{\hr}do\tbr{d}im	&	{\hr}ahdu\tbr{d}um	&	\\
STB	&	\ili{Yamdena}	&	{\hr}ndi\tbr{r}in	&	{\hr}nro\tbr{nr}am	&	{\hr}nru\tbr{nr}um	&	\\
STB	&	\ili{Fordata}	&	{\hr}ri\tbr{d}in	&	{\hr}de\tbr{d}èn	&	{\hr}du\tbr{d}ur	&	\\
STB	&	\ili{Kei}	&	{\hr}ri\tbr{d}in, ri\tbr{n}in	&	{\hr}de\tbr{d}an	&	{\hr}do\tbr{d}	&	\\
STB	&	\ili{Irarutu}	&	\mc{2}{l}{{\hr}ri\tbr{nd}ən, rí\tbr{d}ənə}		&		&	\\
STB	&	\ili{Kowiai}	&	{\hr}ri\tbr{d}in	&		&	{\hr}du\tbr{d}un	&	\\
STB	&	\ili{Geser}	&	{\hr}mari\tbr{h}i	&		&	{\hr}du\tbr{d}un	&	\\
\ili{Banda}	&	\ili{Banda}	&	{\hr}ri\tbr{nd}in	&		&		&	\\
\tsc{AS}	&	\ili{Alune}	&	{\hr}ndi\tbr{r}i	&		&		&	\\
\tsc{AS}	&	\ili{Wakasihu}	&	{\hr}li\tbr{r}inu	&		&		&	\\
\tsc{AS}	&	Kaitetu	&	\mc{2}{l}{{\hr}pa-ri\tbr{k}i (reg. *nd > \it{k})}			&	\\
\tsc{AS}	&	\ili{Saparua}	&	\mc{2}{l}{{\hr}pi-ri\tbr{k}i (reg. *nd > \it{k})}			&	\\
\tsc{AS}	&	\ili{Kayeli}	&	{\hr}namni\tbr{r}i	&		&		&	\\
\tsc{SB}	&	\ili{Buru}	&	{\hr}eb-ri\tbr{d}i	&	{\hr}re\tbr{d}e(n)	&		&	\\
		\lspbottomrule
	\end{tabular}
		\begin{tablenotes}\tnsize
			%\item[†]	{\ili{Kambera} ‹tau ringu› ‘people from cool areas,
								%from the mountainous country’ \citep[442]{Onvlee1984}.}
			\item[†]	{Many of these languages suggest a proto-form like **dumdum ‘thunder, rumble’.}
		\end{tablenotes}
	\end{threeparttable}
%\end{adjustbox}
\end{table}
